% Analysis of the adapter board (front end and references). Macros from results_fe.tex.
\section{Effects of the adapter board}
\label{sec:board}

\subsection{Gain and offset}
The measured transitions go from \SI{\TFirst}{V} to \SI{\TLast}{V} instead of
\SI{\FeTFirstIdeal}{V} to \SI{\FeTLastIdeal}{V}: the LSB is \SI{\Lsb}{mV} instead of
\SI{\FeLsbIdeal}{mV} (gain error \SI{\FeGainErr}{\percent}) and the middle of the range is at
\SI{\FeMid}{V} instead of \SI{\BoardVmid}{V} (\SI{\FeOffsetMv}{mV}, \FeOffsetLsb{} LSB). The
front end doubles any error of the \SI{0.6}{V} of R1/R2 and adds the offsets of the two halves of
U5~\cite{lmv358}; the references are known only from a multimeter reading
(\SI{\BoardVrefn}{V} and \SI{\BoardVrefp}{V}), and \SI{10}{mV} on either is \SI{2.5}{\percent} of the
\SI{0.4}{V} span. Neither error limits the ENOB: both are removed by the
end-point line and by the sine fit.

\subsection{Static linearity is not set by the front end}
At low frequency the front end is linear: the op-amps have all their loop gain, and H2, the
harmonic a single-ended stage would produce, is below \SI{-45}{dBc}. The SINAD measured with
the sines below \SI{50}{kHz} is the same, within \SI{0.3}{dB}, as the SINAD predicted from the
ramp transfer alone (SINAD$_{st}$ in the sweep table), so the low-frequency ENOB of
\SineRefEnob{} bits is entirely the static nonlinearity of the conversion: noise, jitter and
dynamic errors are negligible in comparison. An ideal \AdcBits-bit quantizer would give
$20 \log_{10}(A \sqrt{6}) \approx \SI{30.2}{dB}$ (4.7 bits) for the same amplitude.

The DNL repeats with a period of 8 codes and the INL jumps at codes 8, 16 and 24, the signature
of errors in the binary weights of the capacitive DAC. A least-squares fit of the transition levels
to $T_k = \sum_b w_b\, b_b(k) + c$, with $b_b(k)$ the bits of $k$, gives
\begin{center}\small
\begin{tabular}{lccccc}
\toprule
Bit & 0 & 1 & 2 & 3 & 4 \\
\midrule
Weight (LSB) & \BitWA & \BitWB & \BitWC & \BitWD & \BitWE \\
Ideal & 1 & 2 & 4 & 8 & 16 \\
Error (\%) & \BitEA & \BitEB & \BitEC & \BitED & \BitEE \\
\bottomrule
\end{tabular}
\end{center}
and the INL left after the fit is \InlResidRms{} LSB rms (\InlRms{} LSB before). Errors of several
percent in units made of tens of MOS capacitors are not random mismatch; they come from the chip
(the layout parasitics of the DAC, which the schematic simulation of
Section~\ref{sec:comparison} does not include) or from the references. Because the references
have no decoupling capacitor at the chip (Figure~\ref{fig:refs}), the charge that the DAC draws
at every bit decision (\SI{100}{ns}) is supplied by the trace capacitance and by the open-loop
output impedance of a \SI{1}{MHz} op-amp, which cannot correct it within the bit time. Such a
disturbance depends on the bits already decided, so it would also appear as weight errors.
The measurements of this report cannot separate the two causes. Repeating the ramp with
\SI{100}{nF} and \SI{1}{\micro F} ceramic capacitors at pins 44 and 45 would: if the INL does not
change, the errors are in the chip.

\subsection{Dynamic performance is limited by the front end}
The two outputs of the front end do not have the same path: $v_p$ goes through one op-amp at
unity gain (bandwidth $\approx$ GBW), $v_n$ through that op-amp and then the inverting stage, which
has a noise gain of 2 (bandwidth $\approx$ GBW/2). With the LMV358 (GBW about \SI{1}{MHz}, slew rate
about \SI{1}{V/\micro s}~\cite{lmv358}) both are close to the frequencies of the sweep. A linear
model with single-pole op-amps and the RC network of Figure~\ref{fig:frontend} reproduces the
measured roll-off of the ADC amplitude up to \SI{500}{kHz} with GBW = \SI{\FeGbw}{MHz}
(Figure~\ref{fig:fe}). It has three consequences:
\begin{enumerate}
  \item \textbf{Amplitude loss.} At \SI{\FeHighFreq}{kHz} the differential amplitude falls to
        \FeHighRel{} of its low-frequency value. The ENOB is referred to the amplitude actually
        converted, so a smaller signal alone costs ENOB: the SINAD$_q$ column of
        Table~\ref{tab:fe} is what an ideal quantizer gives for the measured amplitude.
  \item \textbf{Common-mode ripple.} The phase lag of $v_n$ against $v_p$ turns part of the signal
        into common mode at the input frequency: \FeHighCm\,\% of the differential amplitude at
        \SI{\FeHighFreq}{kHz}. With $v_p$ and $v_n$ no longer opposite, the second-order terms of the
        sampling switches and MOS capacitors of the two sides no longer cancel, which could produce H2
        ($v_p^2 - v_n^2 = v_d \cdot 2 v_{cm}$). The measured H2 stays at \SI{-45}{dBc} while the
        ripple is small and rises to \SI{-24}{dBc} once the ripple exceeds about \SI{20}{\percent} of
        the signal, above \SI{200}{kHz} (Figure~\ref{fig:fe}). The transistor-level simulation with the
        modelled ripple (Section~\ref{sec:comparison}) shows, however, that the converter rejects it
        well (H2 below \SI{-50}{dBc}). The measured H2 therefore comes from the op-amps themselves
        near their slew limit (next item), which the linear model does not include; the ripple is a
        symptom of the same bandwidth problem, not the direct cause.
  \item \textbf{Slew rate.} At \SI{\FeHighFreq}{kHz} the input needs \SI{\FeHighSlew}{V/\micro s},
        close to the slew rate of the LMV358, which also has to charge \SI{2.5}{nF} through
        \SI{47.5}{\ohm}. The measured amplitude (\FeHighRel{} of the low-frequency value) is below
        the linear model (\FeHighModel): the op-amps slew, which distorts the signal and explains
        the extra loss.
\end{enumerate}
The circuit of SBOA216~\cite{sboa216} works as intended; its bandwidth is simply that of its
op-amps. For signals up to \SI{650}{kHz} it needs op-amps with tens of MHz of GBW and
\SI{10}{V/\micro s} of slew rate, or a fully differential amplifier, which drives both outputs
through the same path and keeps them opposite.

\begin{figure}[H]\centering
\includegraphics{fig/frontend.pdf}
\caption{Top: differential gain of the front end, measured (fitted ADC amplitude relative to the
lowest frequency) and from the linear model. Middle: common-mode ripple at the ADC input relative to
the differential signal (model). Bottom: measured H2.}
\label{fig:fe}
\end{figure}
\begin{table}[H]\centering\small
\input{table_fe.tex}
\caption{Front-end analysis. $A$: fitted ADC amplitude; diff.\ gain relative to the lowest
frequency; CM/diff.: common-mode ripple of the model; slew: slew rate needed for the input
amplitude; SINAD$_q$: SINAD of an ideal \AdcBits-bit quantizer at the amplitude $A$.}
\label{tab:fe}
\end{table}

\subsection{Conclusions on the measurement setup}
\begin{itemize}
  \item Up to \SI{50}{kHz} the front end does not limit the measurement: the ENOB of
        \SineRefEnob{} bits is the static nonlinearity of the converter, mainly binary weight errors
        of the DAC.
  \item Above \SI{100}{kHz} the front end degrades the measurement: the LMV358 halves of
        the SBOA216 converter have different bandwidths, so the amplitude falls, the input common
        mode moves at the signal frequency (H2) and, near \SI{650}{kHz}, the op-amps slew. The
        drop of ENOB with frequency is mostly the board, not the ADC.
  \item The references have no decoupling at the chip. This may contribute to the weight errors
        and should be fixed before the static linearity is attributed to the chip.
  \item The supply of the ADC is \SI{1.2}{V} and its references are at
        \SI{\BoardVrefn}{V}/\SI{\BoardVrefp}{V}, so the simulations used for the comparison are run
        at these conditions, not at the \SI{1.8}{V} of the design testbench.
\end{itemize}
